% !TeX root = ../../tesis.tex
% Esquema ilustrativo (decisión del autor 2026-11-04): claves y rutas
% tomadas de VFTModel/src/api/main.py; valores = marcadores de tipo.

\section{Ejemplo de Respuesta: Centralidad de Intermediación \texorpdfstring{($B(v)$)}{(B(v))}}
\label{anx:vft-out-betweenness}

Esquema ilustrativo de la respuesta del indicador $B(v)$
(§\ref{subsec:centralidad-intermediacion}). Los nombres de las claves y la ruta son los
de \texttt{VFTModel}; los valores son marcadores de tipo. Las cifras de la
corrida 2026-10-04 se reportan en el Capítulo~\ref{ch:analisis-red-actual}.

El orden va de mayor a menor $B(v)$ y se recorta al valor de \texttt{limit}.

\begin{lstlisting}[style=estiloPython,
  caption={Esquema de respuesta del endpoint
           \texttt{/topological/betweenness-centrality} (ilustrativo).},
  label=lst:out-betweenness]
{
  "status": "success",
  "parametros": { "modo_grafo": "REALISTIC_INTEGRATION",
                  "tolerancia_transbordo_m": <número> },
  "data": {
    "summary": {
      "total_stations_ranked": <número>,
      "top_5": [ { "...": "mismas claves que ranking" } ],
      "reference": { "expected_top_cdmx": ["<estación>"],
                     "expected_B_approx": <número> }
    },
    "ranking": [
      { "node_id": "<texto>", "nombre": "<estación>", "sistema": "<sistema>",
        "alcaldia_municipio": "<texto>", "es_cetram": <booleano>,
        "lon": <número>, "lat": <número>,
        "betweenness_centrality": <número> }
    ]
  }
}
\end{lstlisting}
